This section provides the entry point for the ML calibrated physical
path loss and spread priors and the \textsc{HARP-Net CKM} residual architecture.
The propagation ingredients follow classical FSPL, two ray,
empirical NLoS and large scale parameter models
\cite{rappaport,wocc2021,cost231,tr38901,winner2,khawaja_survey}, while the
residual network inherits dense prediction, FiLM conditioning and
distribution aware regression ideas
\cite{pmnet2023,perez2018film,kendall2017uncertainties,saleh2021probabilistic}.
The detailed mathematical derivations and implementation level fitting notes
are included in this methodology chapter, so the final pipeline can be read
without leaving the method section.

\subsection{High level architecture}

The central design decision in the final model is \emph{not} to learn path
loss, delay spread, and angular spread from scratch. Instead:

\begin{itemize}
\setlength{\itemsep}{0.12em}
\setlength{\parsep}{0pt}
\setlength{\topsep}{0.15em}
\item the path loss prior provides a frozen prior for \texttt{path\_loss},
\item the spread prior module provides frozen priors for delay spread
      (\texttt{delay\_spread}) and angular spread
      (\texttt{angular\_spread}),
\item \textsc{HARP-Net CKM} learns bounded corrections around those priors.
\end{itemize}

The deterministic or interpretable structure is kept in the prior, and the deep
model is used only for the residual part that the prior cannot explain.
Figure~\ref{fig:methodology_full_framework} illustrates the overall pipeline.

This prior phase is not presented as a new universal closed form propagation
law. Its value is more practical: it turns physics, morphology, visibility, and
height dependent calibration into a strong auditable baseline. The final model
then follows a split of responsibility: calibrated priors
remove the stable component of the targets, and the neural residual model
focuses on the local errors that remain.

The term \emph{prior} is not used here in the Bayesian sense of a probability
distribution over model parameters or targets, and no posterior inference is
performed in the prior stage. It denotes deterministic, frozen estimator maps
computed before the neural residual correction. The word \emph{prior} is
therefore used for three concrete estimators, not for a vague physical hint:
\begin{itemize}
\setlength{\itemsep}{0.16em}
\setlength{\parsep}{0pt}
\setlength{\topsep}{0.15em}
\item \textbf{LoS path loss prior.} This is the clearest physical branch:
      \(\mathrm{PL}_{\mathrm{LoS}}=\mathrm{FSPL}+C_{\mathrm{2ray}}+b+r\).
      FSPL uses the 3D direct Tx--Rx distance; \(C_{\mathrm{2ray}}\) comes
      from a coherent direct plus ground reflected field sum
      \cite{rappaport,wocc2021}; \(b\) is a fitted height dependent dB offset;
      and \(r\) is a smoothed radial
      residual table. The UAV is always centred horizontally, so the radial
      coordinate is unambiguous, but \(h_{\mathrm{tx}}\) varies by sample:
      \(\rho(h_{\mathrm{tx}})\), \(\phi(h_{\mathrm{tx}})\), \(b(h_{\mathrm{tx}})\)
      and \(r(d_{2D},h_{\mathrm{tx}})\) are fitted in \SI{5}{m} height bins,
      smoothed, and linearly interpolated at inference. It is not a generic
      free space formula and it is not a learned linear regressor.
\item \textbf{NLoS path loss prior.} This starts from the raw blocked link
      propagation estimate \(\mathrm{PL}^{(0)}\), which includes
      \(\mathrm{PL}_{\mathrm{NLoS,raw}}\) on obstructed pixels, then calibrates
      it with local morphology: building density, building height, NLoS support,
      the elevation dependent shadowing proxy \(\sigma_{\mathrm{shadow}}\),
      elevation angle, and interactions. It uses antenna height regimes
      (\texttt{low\_ant}, \texttt{mid\_ant}, \texttt{high\_ant}) because the
      obstruction geometry changes qualitatively when the centred UAV rises
      above the local skyline, consistent with A2G elevation and scenario
      modelling practice \cite{cost231,tr38901,alhourani2014,khawaja_survey}.
      It is deliberately separate from the LoS branch
      because blocked pixels do not follow the same two ray mechanism.
\item \textbf{Delay and angular spread priors.} These are fitted in
      \(\log(1+y)\) space. The raw estimate captures propagation state,
      distance, elevation, and local morphology; the final ridge calibration
      adds the 23-feature vector of Sec.~\ref{sec:try79_detail}. That vector
      includes continuous normalized UAV height, density weighted transmitter
      clearance over nearby building pixels, and the local window area fraction
      occupied by buildings taller than the transmitter. This follows the
      log domain convention used for large scale spread statistics and recent
      A2G spread predictors \cite{tr38901,winner2,yang2019a2gml,huang2025a2gtransformer}.
      This is why the spread prior can fit the near zero LoS mass
      and the NLoS tail without letting the largest tail values dominate the
      whole calibration, unlike a native space global average.
\end{itemize}

The notation intentionally reuses \(\rho\) in two different blocks:
\(\rho(h_{\mathrm{tx}})\) is the fitted LoS reflection amplitude, whereas
\(\rho_k\) is a local building density fraction in a \(k\times k\) window.
They are unrelated quantities; the argument or subscript identifies which one
is meant.
The same convention is used for \(b\): \(b(h_{\mathrm{tx}})\) is a LoS
dB bias, \(b(x)\) is the binary building mask, \(b_{\mathrm{topology}}\) is
a spread prior log domain offset, and \(b_{41}\) is the blocker depth feature.
These symbols are local to their definitions and should not be read as one
shared physical variable.

\begin{table}[H]
\centering
\caption{Term by term interpretation of the main frozen prior quantities.}
\label{tab:prior_term_meanings}
\scriptsize
\renewcommand{\arraystretch}{1.0}
\begin{tabularx}{\textwidth}{>{\raggedright\arraybackslash}p{0.18\textwidth}
>{\raggedright\arraybackslash}p{0.20\textwidth}
>{\raggedright\arraybackslash}X}
\toprule
Term & Prior block & Interpretation \\
\midrule
\(d_{\mathrm{LoS}}\) & LoS path loss &
3D direct Tx--Rx distance. It sets the FSPL baseline and uses the fixed
\(h_{\mathrm{rx}}=\SI{1.5}{m}\) receiver height. \\
\(d_{\mathrm{ref}}\) & LoS path loss &
Image source reflected distance. It is the length of the single
ground reflected path used in the coherent two ray correction. \\
\(C_{\mathrm{2ray}}\) & LoS path loss &
Direct plus reflected field correction in dB. Negative values mean constructive
gain over FSPL; positive values mean destructive fading. \\
\(\rho(h_{\mathrm{tx}})\) & LoS path loss &
Height interpolated effective reflection amplitude. It is fitted from CKM
rather than assumed from one fixed ground material. \\
\(\phi(h_{\mathrm{tx}})\) & LoS path loss &
Height interpolated effective reflection phase. It absorbs the phase offset of
the simplified one reflection geometry. \\
\(b(h_{\mathrm{tx}})\) & LoS path loss &
Height dependent dB bias remaining after the coherent two ray fit. \\
\(r(d_{2D},h_{\mathrm{tx}})\) & LoS path loss &
Smoothed radial residual table. It corrects distance dependent errors that
remain after FSPL, two ray interference, and height bias. \\
\(\mathrm{PL}^{(0)}\), \(\mathrm{PL}_{\mathrm{NLoS,raw}}\) & NLoS path loss &
Raw hybrid formula map and its obstructed pixel branch before
map morphology calibration. \\
\(\rho_k\) & NLoS PL and spreads &
Fraction of building pixels in a \(k\times k\) neighbourhood; a local density
proxy. \\
\(h_k\) & NLoS PL and spreads &
Local building height summary in a \(k\times k\) neighbourhood; a proxy for
nearby reflectors and blockers. \\
\(n_k\) & NLoS PL and spreads &
Local NLoS support fraction; indicates how much of the neighbourhood is
geometrically blocked. \\
\(\sigma_{\mathrm{shadow}}\) & NLoS path loss &
Elevation dependent shadowing proxy used by the NLoS linear calibration. \\
\(\theta\), \(\theta_{\mathrm{norm}}\), \(\theta_{\mathrm{inv}}\) &
PL and spreads &
Elevation angle terms. Low elevation links interact with more buildings, so
\(\theta_{\mathrm{inv}}\) is especially useful for spread tails. \\
\(b_{41}\), \(c_{41}\), \(t_{41}\) & Spreads &
Blocker depth, transmitter clearance, and local taller building fraction; they
separate open overhead links from obstructed street canyon like links. \\
\bottomrule
\end{tabularx}
\end{table}

\subsection{Prior design rules}
\label{sec:good_prior_design}

The main prior design rule is that a useful prior is not just a physical
looking equation. It is a complete estimator: the formula, fitted calibration,
feature units, LoS/NLoS mask, topology features, and evaluation mask must be
identical between fitting and inference. The final prior is strong because it
keeps those pieces self contained and auditable under the same city holdout
contract used for the final model.

The distinction matters most for NLoS. A calibration fitted for one physical
formula cannot reliably correct another; local height and density features
must have the same units when fitted and when applied; and regime assignment
should come from the sample geometry rather than only from a broad city label.
Once those constraints were enforced, the prior could separate deterministic
large scale structure from the genuinely local residual that the final model later
learned.

This also explains why the final network is residual and distribution aware.
Blocked NLoS and spread targets are multimodal and heavy tailed, so a single
point estimate without a strong anchor tends to average incompatible regimes.
\textsc{HARP-Net CKM} therefore keeps the priors as frozen anchors and lets the
neural head model bounded residual structure and uncertainty without turning
the main method into an unconstrained direct regressor.

\begin{table}[H]
\centering
\caption{Compact design rules extracted from the prior experiments.}
\label{tab:prior_design_rules}
\small
\begin{tabularx}{\textwidth}{>{\raggedright\arraybackslash}p{0.28\textwidth}
>{\raggedright\arraybackslash}X}
\toprule
Rule & Effect on the final pipeline \\
\midrule
Fit and apply the same formula &
The calibration corrects the actual runtime prior, not a neighbouring variant. \\
Keep units explicit &
Distances, heights, topology values, and normalized features have one shared
definition across fitting, validation, and test. \\
Use the metric mask during fitting &
The prior is optimized for valid receiver pixels rather than for building cells
that are excluded from the final metric. \\
Hard separate LoS and NLoS &
The two regimes keep different physics, feature sets, and calibration
parameters when the dataset provides a reliable mask. \\
Let literature define axes, not all coefficients &
Distance, height, frequency, and elevation come from propagation theory;
effective reflection and local morphology corrections are fitted from CKM. \\
Learn bounded residuals after the prior &
The final model is allowed to correct local errors without overwriting the calibrated
large scale structure. \\
Treat calibration artifacts as part of the model &
Large JSON and coefficient artifacts are not bookkeeping only; they store the
height, density, visibility, and morphology-specific decisions that make the
prior a competitive baseline. \\
\bottomrule
\end{tabularx}
\end{table}

\subsection{Literature ingredients and thesis adaptations}

The final pipeline is a hybrid between classical propagation models,
distribution aware regression, and dataset-specific calibration.
Table~\ref{tab:prior_literature_mapping} keeps the mapping explicit; the
detailed equations and fitting procedures follow in this chapter, covering the
attenuation prior, spread prior, and \textsc{HARP-Net CKM} residual model.

\begin{table}[H]
\centering
\caption{Literature to design mapping used by \textsc{HARP-Net CKM}.}
\label{tab:prior_literature_mapping}
\scriptsize
\renewcommand{\arraystretch}{1.12}
\begin{tabularx}{\textwidth}{>{\raggedright\arraybackslash}p{0.18\textwidth}
>{\raggedright\arraybackslash}X
>{\raggedright\arraybackslash}X
>{\raggedright\arraybackslash}X}
\toprule
Block & Literature basis & CKM adaptation & New role in this thesis \\
\midrule
LoS path loss prior &
FSPL and coherent two ray propagation \cite{rappaport,wocc2021}. &
Effective complex reflection, height bin interpolation, and radial residual
fitting on CKM LoS pixels. &
Reconstructs the radial interference pattern and supplies the frozen LoS
anchor. \\
NLoS path loss calibration &
COST-231/Hata style attenuation, A2G elevation envelopes, and LoS/NLoS
separation \cite{cost231,tr38901,alhourani2014,khawaja_survey}. &
Map derived density, height, NLoS support, elevation, and shadowing features
with regime wise calibration. &
Turns a broad empirical curve into a local, topology aware NLoS prior. \\
Delay and angular spread priors &
Log normal spread statistics from 3GPP/WINNER style channel models and UAV
elevation trends \cite{tr38901,winner2,yang2019a2gml,huang2025a2gtransformer}. &
Log domain ridge regression with local morphology features and a sparse regime
fallback hierarchy. &
Provides stable spread anchors without pretending that spreads are smooth
deterministic channel attenuation maps. \\
Final residual model &
U-Net style dense prediction, FiLM conditioning, probabilistic regression, and
mixture density ideas
\cite{isola2017pix2pix,pmnet2023,perez2018film,kendall2017uncertainties,saleh2021probabilistic,garciamarti2020mixture}. &
Prior channels, explicit masks, residual caps, anchor gates, and a shared
multitask trunk. &
Learns only the bounded local correction and uncertainty left after the
calibrated priors. \\
\bottomrule
\end{tabularx}
\end{table}
